% !TeX program = lualatex
% !TeX encoding = utf8
% !TeX spellcheck = uk_UA
% !TeX root = ../main.tex


%=========================================================
\chapter[Рух заряджених частинок]{Рух заряджених частинок в~електричних і~магнітних полях}\label{\currfilebase}\hypertarget{EBMotion}{}
\graphicspath{{\currfilebase/Pictures}}
%=========================================================
\epigraph{\Alexander Я не бачу способу уникнути висновку, що це~--- заряди негативної електрики, які переносяться частинками речовини. Постає
	питання: що це за частинки~--- атоми, молекули чи матерія в~ще тоншому стані поділу?
}{Дж.~Дж.~Томсон, \enquote{Cathode Rays}, 1897}


%% --------------------------------------------------------
\section{Рух частинок в~однорідних полях}\hypertarget{motion-in-uniform-fields}{}
%% --------------------------------------------------------


%% --------------------------------------------------------
\subsection*{Рух в~однорідному електричному полі}
%% --------------------------------------------------------


Якщо напруженість поля $\Efield = \const$, то з~рівняння руху
% ~~~~~~~~~~~~~~~~~~~~~~~~~~~~~~~~~~~~~~~~~~~~~~~~~~~~~~~~
\begin{equation}\label{eq:motion_E}
	m\dot{\vect{v}} = q\Efield
\end{equation}
% ~~~~~~~~~~~~~~~~~~~~~~~~~~~~~~~~~~~~~~~~~~~~~~~~~~~~~~~~
випливає
% ~~~~~~~~~~~~~~~~~~~~~~~~~~~~~~~~~~~~~~~~~~~~~~~~~~~~~~~~
\begin{equation}\label{eq:motion_E_solution}
	\vect{v} = \vect{v}_0 + \frac{q\Efield}{m}t, \qquad
	\vect{r} = \vect{r}_0 + \vect{v}_0 t + \frac{q\Efield}{2m}t^2,
\end{equation}
% ~~~~~~~~~~~~~~~~~~~~~~~~~~~~~~~~~~~~~~~~~~~~~~~~~~~~~~~~
тобто маємо рівноприскорений рух, де прискорення напрямлене уздовж вектора напруженості поля.


%% --------------------------------------------------------
\subsection*{Рух в~однорідному магнітному полі}
%% --------------------------------------------------------


На частинку діє сила Лоренца, так що рівняння руху має вигляд
% ~~~~~~~~~~~~~~~~~~~~~~~~~~~~~~~~~~~~~~~~~~~~~~~~~~~~~~~~
\begin{equation}\label{eq:Lorentz_motion}
	m\dot{\vect{v}} = \frac{q}{c}\vect{v} \times \vect{B}.
\end{equation}
% ~~~~~~~~~~~~~~~~~~~~~~~~~~~~~~~~~~~~~~~~~~~~~~~~~~~~~~~~
Очевидно, що $\dot{\vect{v}} \perp \vect{v}$, $\dot{\vect{v}} \perp \vect{B}$. Перша рівність означає, що сила Лоренца не виконує роботи: $\dfrac{d}{dt}\left(\dfrac{mv^2}{2}\right) = m\vect{v}\cdot\dot{\vect{v}} = 0$, тобто $|\vect{v}| = \const$~--- магнітне поле змінює лише напрямок швидкості. Розкладемо вектор швидкості на складові, паралельну і~перпендикулярну полю: $\vect{v} = \vect{v}_{\parallel} + \vect{v}_{\perp}$. Для цих складових маємо рівняння
% ~~~~~~~~~~~~~~~~~~~~~~~~~~~~~~~~~~~~~~~~~~~~~~~~~~~~~~~~
\begin{equation}\label{eq:v_parallel_motion}
	m\dot{\vect{v}}_{\parallel} = 0,
\end{equation}
% ~~~~~~~~~~~~~~~~~~~~~~~~~~~~~~~~~~~~~~~~~~~~~~~~~~~~~~~~
% ~~~~~~~~~~~~~~~~~~~~~~~~~~~~~~~~~~~~~~~~~~~~~~~~~~~~~~~~
\begin{equation}\label{eq:v_perp_motion}
	m\dot{\vect{v}}_{\perp} = \frac{q}{c}\vect{v}_{\perp} \times \vect{B}.
\end{equation}
% ~~~~~~~~~~~~~~~~~~~~~~~~~~~~~~~~~~~~~~~~~~~~~~~~~~~~~~~~
З~рівняння \eqref{eq:v_parallel_motion} випливає $\vect{v}_{\parallel} = \const$.

Друге рівняння перепишемо у~вигляді
% ~~~~~~~~~~~~~~~~~~~~~~~~~~~~~~~~~~~~~~~~~~~~~~~~~~~~~~~~
\begin{equation}\label{eq:v_perp_rotation}
	\dot{\vect{v}}_{\perp} = \vect{\omega} \times \vect{v}_{\perp}, \qquad
	\vect{\omega} = -\omega_c\hat{\vect{b}}, \qquad
	\omega_c = \frac{qB}{mc}, \quad \hat{\vect{b}} = \frac{\vect{B}}{B}.
\end{equation}
% ~~~~~~~~~~~~~~~~~~~~~~~~~~~~~~~~~~~~~~~~~~~~~~~~~~~~~~~~

Воно описує обертання навколо напрямку магнітного поля з~кутовою швидкістю $|\omega_c| = |q|B/mc$, яка називається \emphx{циклотронною частотою},~--- рис.~\ref{tikz:cyclotron}. Величина $\omega_c$ є знаковою: для $q<0$ вона від'ємна, і~тоді напрямок обертання змінюється на протилежний.

\begin{figure}[h!]
	\centering
	\localinput{CycletronicTrajectory.tikz}
	\caption{Обертальний рух частинки із зарядом $q > 0$ навколо напрямку магнітного поля.}
	\label{tikz:cyclotron}
\end{figure}

Якщо в~момент потрапляння в~область дії поля складова швидкості частинки $\vect{v}_{\perp}$ дорівнювала (за величиною) $v_{\perp}$, то радіус кола дорівнює
% ~~~~~~~~~~~~~~~~~~~~~~~~~~~~~~~~~~~~~~~~~~~~~~~~~~~~~~~~
\begin{equation}\label{eq:cyclotron_radius}
	R = \frac{v_{\perp}}{|\omega_c|} = \frac{mcv_{\perp}}{|q|B}.
\end{equation}
% ~~~~~~~~~~~~~~~~~~~~~~~~~~~~~~~~~~~~~~~~~~~~~~~~~~~~~~~~
Величину $R$ називають \emphx{ларморівським радіусом}. Якщо $R$ набагато менший за характерний розмір області, де діє поле, то частинка встигає багато разів обернутися навколо силової лінії, і~її рух упоперек поля сильно обмежений; такі частинки називають \emphx{замагніченими}.


%% --------------------------------------------------------
\subsection*{Рух в~однорідних паралельних електричному та магнітному полях ($\Efield \parallel \vect{B}$)}
%% --------------------------------------------------------


У~рівнянні руху
% ~~~~~~~~~~~~~~~~~~~~~~~~~~~~~~~~~~~~~~~~~~~~~~~~~~~~~~~~
\begin{equation}\label{eq:motion_EB_parallel}
	m\dot{\vect{v}} = q\Efield + \frac{q}{c}\vect{v} \times \vect{B}
\end{equation}
% ~~~~~~~~~~~~~~~~~~~~~~~~~~~~~~~~~~~~~~~~~~~~~~~~~~~~~~~~
покладемо $\vect{v} = \vect{v}_{\parallel} + \vect{v}_{\perp}$, де компоненти швидкості відповідно $\vect{v}_{\parallel} \parallel \vect{B}$ і~$\vect{v}_{\perp} \perp \vect{B}$. Для цих компонент отримуємо рівняння
% ~~~~~~~~~~~~~~~~~~~~~~~~~~~~~~~~~~~~~~~~~~~~~~~~~~~~~~~~
\begin{equation}\label{eq:v_parallel_E}
	m\dot{\vect{v}}_{\parallel} = q\Efield,
\end{equation}
% ~~~~~~~~~~~~~~~~~~~~~~~~~~~~~~~~~~~~~~~~~~~~~~~~~~~~~~~~
% ~~~~~~~~~~~~~~~~~~~~~~~~~~~~~~~~~~~~~~~~~~~~~~~~~~~~~~~~
\begin{equation}\label{eq:v_perp_B}
	m\dot{\vect{v}}_{\perp} = \frac{q}{c}\vect{v}_{\perp} \times \vect{B}.
\end{equation}
% ~~~~~~~~~~~~~~~~~~~~~~~~~~~~~~~~~~~~~~~~~~~~~~~~~~~~~~~~

\begin{SCfigure}[2][h!]
	\centering
	\localinput{orbit.tikz}
	\caption{Траєкторія частинки в~паралельних електричному і~магнітному полях.}
	\label{tikz:helix}
\end{SCfigure}

Перше з~них описує рівноприскорений рух уздовж напрямку електричного поля, а~друге рівняння~--- обертальний рух навколо напрямку магнітного поля з~циклотронною частотою $|\omega_c| = |q|B/mc$. Траєкторія частинки являє собою спіраль зі зростаючим кроком і~з~віссю, паралельною силовій лінії магнітного поля (рис.~\ref{tikz:helix}).


%% --------------------------------------------------------
\subsection*{Рух у~схрещених полях ($\Efield \perp \vect{B}$)}
%% --------------------------------------------------------


Як і~вище, виділимо компоненти швидкості, паралельну та перпендикулярну вектору магнітного поля: $\vect{v} = \vect{v}_{\parallel} + \vect{v}_{\perp}$ (рис.~\ref{tikz:drift}). Для цих компонент маємо рівняння
% ~~~~~~~~~~~~~~~~~~~~~~~~~~~~~~~~~~~~~~~~~~~~~~~~~~~~~~~~
\begin{equation}\label{eq:v_parallel_crossed}
	m\dot{\vect{v}}_{\parallel} = 0,
\end{equation}
% ~~~~~~~~~~~~~~~~~~~~~~~~~~~~~~~~~~~~~~~~~~~~~~~~~~~~~~~~
% ~~~~~~~~~~~~~~~~~~~~~~~~~~~~~~~~~~~~~~~~~~~~~~~~~~~~~~~~
\begin{equation}\label{eq:v_perp_crossed}
	m\dot{\vect{v}}_{\perp} = q\left(\Efield + \frac{1}{c}\vect{v}_{\perp} \times \vect{B}\right).
\end{equation}
% ~~~~~~~~~~~~~~~~~~~~~~~~~~~~~~~~~~~~~~~~~~~~~~~~~~~~~~~~

Рівняння \eqref{eq:v_parallel_crossed} дає $\vect{v}_{\parallel} = \const$, що означає рух з~постійною швидкістю вздовж вектора $\vect{B}$, яка визначається початковими умовами.

\begin{figure}[h!]
	\centering
	\localinput{eldriftplot.tikz}
	\caption{Електричний дрейф додатно (синя крива) та від'ємно (червона крива) заряджених частинок.}
	\label{tikz:drift}
\end{figure}

Для знаходження компоненти $\vect{v}_{\perp}$ покладемо $\vect{v}_{\perp} = \vect{u} + \vect{v}'$ і~підберемо вектор $\vect{u}$ так, щоб виключити електричне поле:
% ~~~~~~~~~~~~~~~~~~~~~~~~~~~~~~~~~~~~~~~~~~~~~~~~~~~~~~~~
\begin{equation}\label{eq:u_condition}
	\Efield + \frac{1}{c}\vect{u} \times \vect{B} = 0.
\end{equation}
% ~~~~~~~~~~~~~~~~~~~~~~~~~~~~~~~~~~~~~~~~~~~~~~~~~~~~~~~~
Оскільки $\vect{u}\times\vect{B} \perp \vect{B}$, це рівняння має розв'язок лише за умови $\Efield \perp \vect{B}$ (саме такий випадок розглядаємо). Складова $\vect{u}$, паралельна $\vect{B}$, ні на що не впливає (її можна віднести до $\vect{v}_{\parallel}$), тому вважаємо $\vect{u} \perp \vect{B}$, тобто $\vect{u}\cdot\vect{B} = 0$. Для знаходження вектора $\vect{u}$ помножимо почленно рівняння \eqref{eq:u_condition} векторно зліва на $\vect{B}$:
% ~~~~~~~~~~~~~~~~~~~~~~~~~~~~~~~~~~~~~~~~~~~~~~~~~~~~~~~~
\begin{equation}\label{eq:u_cross_B}
	\vect{B} \times \Efield + \frac{1}{c}\vect{B} \times (\vect{u} \times \vect{B}) = 0.
\end{equation}
% ~~~~~~~~~~~~~~~~~~~~~~~~~~~~~~~~~~~~~~~~~~~~~~~~~~~~~~~~
Перетворимо подвійний векторний добуток з~урахуванням взаємної ортогональності векторів $\vect{B}$ і~$\vect{u}$:
% ~~~~~~~~~~~~~~~~~~~~~~~~~~~~~~~~~~~~~~~~~~~~~~~~~~~~~~~~
\begin{equation}\label{eq:double_cross}
	\vect{B} \times (\vect{u} \times \vect{B}) = \vect{u}B^2 - \vect{B}(\vect{u}\cdot\vect{B}) = \vect{u}B^2.
\end{equation}
% ~~~~~~~~~~~~~~~~~~~~~~~~~~~~~~~~~~~~~~~~~~~~~~~~~~~~~~~~
У~результаті отримаємо
% ~~~~~~~~~~~~~~~~~~~~~~~~~~~~~~~~~~~~~~~~~~~~~~~~~~~~~~~~
\begin{equation}\label{eq:drift_velocity}
	\tcbhighmath{
		\vect{u} = c\frac{\Efield \times \vect{B}}{B^2}.
	}
\end{equation}
% ~~~~~~~~~~~~~~~~~~~~~~~~~~~~~~~~~~~~~~~~~~~~~~~~~~~~~~~~
Цей рух являє собою дрейф (з~постійною швидкістю) в~напрямку, перпендикулярному обом полям ($\Efield$ і~$\vect{B}$) (рис.~\ref{tikz:drift}). Відповідно для компоненти $\vect{v}'$ отримаємо рівняння
% ~~~~~~~~~~~~~~~~~~~~~~~~~~~~~~~~~~~~~~~~~~~~~~~~~~~~~~~~
\begin{equation}\label{eq:v_prime_rotation}
	m\dot{\vect{v}}' = \frac{q}{c}\vect{v}' \times \vect{B},
\end{equation}
% ~~~~~~~~~~~~~~~~~~~~~~~~~~~~~~~~~~~~~~~~~~~~~~~~~~~~~~~~
яке, як і~рівняння \eqref{eq:v_perp_rotation}, описує обертальний рух навколо напрямку магнітного поля з~циклотронною частотою $|\omega_c| = |q|B/mc$.

Сумарний рух являє собою накладення:
\begin{enumerate}
	\item рівномірного руху вздовж вектора $\vect{B}$ зі швидкістю $(\vect{v}_{\parallel})_0$, яка дорівнює компоненті швидкості $\vect{v}_{\parallel}$ в~початковий момент,
	\item дрейфового руху зі швидкістю $\vect{u}$,
	\item обертання навколо напрямку $\vect{B}$.
\end{enumerate}

Дрейфова швидкість $u = cE/B$ можлива лише при $E < B$ (тоді $u < c$). При $E > B$ системи відліку, в~якій електричне поле зникає, не існує, і~частинка необмежено прискорюється. Викладений нерелятивістський розв'язок справедливий при $E \ll B$ (і~$|\vect{v}| \ll c$). Якщо ж $E$ порівнянне з~$B$, слід розв'язувати релятивістські рівняння руху: швидкість дрейфу $cE/B$ при $E<B$ зберігається, але рух у~системі відліку, що дрейфує, вже не є простим обертанням з~частотою $\omega_c$.

Знайдемо траєкторію шляхом прямого розв'язання рівнянь \eqref{eq:v_perp_crossed}. Для цього запишемо ці рівняння в~системі координат, показаній на рис.~\ref{tikz:drift}. Введемо орти $\{\vect{i}, \vect{j}, \vect{k}\}$ відповідно вздовж координатних осей $\{x, y, z\}$. Оскільки в~розглянутому випадку $\vect{v} \times \vect{B} = \vect{i}B_z v_y - \vect{j}B_z v_x$, то, позначаючи для стислості $E_y = E$, $B_z = B$, маємо:
% ~~~~~~~~~~~~~~~~~~~~~~~~~~~~~~~~~~~~~~~~~~~~~~~~~~~~~~~~
\begin{equation}\label{eq:system_components}
	\text{(а) } \dot{v}_z = 0; \qquad
	\text{(б) } \dot{v}_x = \frac{qB}{mc}v_y; \qquad
	\text{(в) } \dot{v}_y = \frac{q}{m}E - \frac{qB}{mc}v_x.
\end{equation}
% ~~~~~~~~~~~~~~~~~~~~~~~~~~~~~~~~~~~~~~~~~~~~~~~~~~~~~~~~

Рівняння \eqref{eq:system_components}а~описує рівномірний рух уздовж вектора магнітного поля:
% ~~~~~~~~~~~~~~~~~~~~~~~~~~~~~~~~~~~~~~~~~~~~~~~~~~~~~~~~
\begin{equation}\label{eq:vz_motion}
	v_z = v_{\parallel} = \const, \qquad z = z_0 + v_{\parallel}t.
\end{equation}
% ~~~~~~~~~~~~~~~~~~~~~~~~~~~~~~~~~~~~~~~~~~~~~~~~~~~~~~~~
Далі вважаємо $v_{\parallel} = 0$. Для розв'язання рівнянь \eqref{eq:system_components}б і~\eqref{eq:system_components}в~покладемо
% ~~~~~~~~~~~~~~~~~~~~~~~~~~~~~~~~~~~~~~~~~~~~~~~~~~~~~~~~
\begin{equation}\label{eq:Vx_substitution}
	V_x = v_x - u, \qquad u = cE/B.
\end{equation}
% ~~~~~~~~~~~~~~~~~~~~~~~~~~~~~~~~~~~~~~~~~~~~~~~~~~~~~~~~
Після цього рівняння \eqref{eq:system_components}б, в~набирають вигляду
% ~~~~~~~~~~~~~~~~~~~~~~~~~~~~~~~~~~~~~~~~~~~~~~~~~~~~~~~~
\begin{equation}\label{eq:Vx_vy_system}
	\dot{V}_x = \omega_c v_y, \qquad \dot{v}_y = -\omega_c V_x,
\end{equation}
% ~~~~~~~~~~~~~~~~~~~~~~~~~~~~~~~~~~~~~~~~~~~~~~~~~~~~~~~~
де $\omega_c = qB/mc$~--- циклотронна частота (для $q<0$ вона від'ємна). Виключаючи звідси змінну $V_x$, зведемо цю систему до рівняння гармонічного осцилятора:
% ~~~~~~~~~~~~~~~~~~~~~~~~~~~~~~~~~~~~~~~~~~~~~~~~~~~~~~~~
\begin{equation}\label{eq:harmonic_oscillator}
	\ddot{v}_y = -\omega_c^2 v_y,
\end{equation}
% ~~~~~~~~~~~~~~~~~~~~~~~~~~~~~~~~~~~~~~~~~~~~~~~~~~~~~~~~
звідки знаходимо:
% ~~~~~~~~~~~~~~~~~~~~~~~~~~~~~~~~~~~~~~~~~~~~~~~~~~~~~~~~
\begin{equation}\label{eq:vy_solution}
	v_y = v_0 \cos(\omega_c t + \varphi_0), \qquad
	y = y_0 + \frac{v_0}{\omega_c} \sin(\omega_c t + \varphi_0).
\end{equation}
% ~~~~~~~~~~~~~~~~~~~~~~~~~~~~~~~~~~~~~~~~~~~~~~~~~~~~~~~~
За допомогою другого рівняння системи \eqref{eq:Vx_vy_system} і~підстановки \eqref{eq:Vx_substitution} знаходимо $v_x(t)$ і~$x(t)$:
% ~~~~~~~~~~~~~~~~~~~~~~~~~~~~~~~~~~~~~~~~~~~~~~~~~~~~~~~~
\begin{equation}\label{eq:vx_x_solution}
	\begin{gathered}
		v_x = u - \frac{1}{\omega_c}\dot{v}_y = u + v_0 \sin(\omega_c t + \varphi_0), \\
		x = x_0 + ut - \frac{v_0}{\omega_c} \cos(\omega_c t + \varphi_0).
	\end{gathered}
\end{equation}
% ~~~~~~~~~~~~~~~~~~~~~~~~~~~~~~~~~~~~~~~~~~~~~~~~~~~~~~~~
Формули \eqref{eq:vy_solution}, \eqref{eq:vx_x_solution} визначають траєкторію частинки за будь-яких початкових умов. Виберемо початок координат так, що $x(0) = 0$, $y(0) = 0$. Тоді
% ~~~~~~~~~~~~~~~~~~~~~~~~~~~~~~~~~~~~~~~~~~~~~~~~~~~~~~~~
\begin{equation}\label{eq:trochoid}
	\begin{cases}
		x = ut + \frac{v_0}{\omega_c}[\cos\varphi_0 - \cos(\omega_c t + \varphi_0)], \\
		y = \frac{v_0}{\omega_c}[\sin(\omega_c t + \varphi_0) - \sin\varphi_0],
	\end{cases}
	\qquad
	\begin{cases}
		v_x = u + v_0 \sin(\omega_c t + \varphi_0), \\
		v_y = v_0 \cos(\omega_c t + \varphi_0).
	\end{cases}
\end{equation}
% ~~~~~~~~~~~~~~~~~~~~~~~~~~~~~~~~~~~~~~~~~~~~~~~~~~~~~~~~
Крива, що задається рівняннями \eqref{eq:trochoid}, називається \emphx{трохоїдою}. Її вигляд залежить від відношення $|v_0|/|u|$, яке визначається початковими умовами: при $|v_0| < |u|$ це хвиляста крива без петель, при $|v_0| > |u|$~--- крива з~петлями, а~при $|v_0| = |u|$~--- циклоїда з~вістрями, в~яких швидкість частинки перетворюється на нуль. Розглянемо частинний випадок, коли початкова швидкість частинки дорівнює нулю:
% ~~~~~~~~~~~~~~~~~~~~~~~~~~~~~~~~~~~~~~~~~~~~~~~~~~~~~~~~
\begin{equation}\label{eq:zero_initial_velocity}
	v_x(0) = v_y(0) = 0.
\end{equation}
% ~~~~~~~~~~~~~~~~~~~~~~~~~~~~~~~~~~~~~~~~~~~~~~~~~~~~~~~~
Тоді $\varphi_0 = \pi/2$, $v_0 = -u$ (тут $v_0$~--- алгебраїчна амплітуда, яка може бути від'ємною), і~з~\eqref{eq:trochoid} знаходимо:
% ~~~~~~~~~~~~~~~~~~~~~~~~~~~~~~~~~~~~~~~~~~~~~~~~~~~~~~~~
\begin{equation}\label{eq:cycloid}
	\begin{gathered}
		v_x = u(1 - \cos\omega_c t), \qquad v_y = u \sin\omega_c t; \\
		x(t) = ut - \frac{u}{\omega_c}\sin\omega_c t, \qquad
		y(t) = \frac{u}{\omega_c}(1 - \cos\omega_c t).
	\end{gathered}
\end{equation}
% ~~~~~~~~~~~~~~~~~~~~~~~~~~~~~~~~~~~~~~~~~~~~~~~~~~~~~~~~

Ця траєкторія являє собою \emphx{циклоїду} (рис.~\ref{tikz:drift_cycloide}). Таку траєкторію описує точка на ободі колеса радіуса $|u/\omega_c|$, що котиться без проковзування по площині зі швидкістю $u$. Відповідно до сказаного траєкторія являє собою суперпозицію дрейфу вздовж осі $x$ і~обертання навколо осі $z$ з~циклотронною частотою $\omega_c = qB/mc$.

\begin{figure}[h!]
	\centering
	\localinput{eldriftplot_cycloide.tikz}
	\caption{Траєкторія частинки у~взаємно перпендикулярних полях $\vect{B}$ і~$\Efield$ за умови, що початкова швидкість дорівнює нулю.}
	\label{tikz:drift_cycloide}
\end{figure}


%% --------------------------------------------------------
\subsection*{Загальний випадок схрещених полів}
%% --------------------------------------------------------


У~загальному випадку є складова вектора електричного поля, паралельна магнітному полю: $\Efield_{\parallel} \neq 0$. Тоді для компоненти швидкості $\vect{v}_{\parallel}$ маємо рівняння
% ~~~~~~~~~~~~~~~~~~~~~~~~~~~~~~~~~~~~~~~~~~~~~~~~~~~~~~~~
\begin{equation}\label{eq:v_parallel_general}
	m\dot{\vect{v}}_{\parallel} = q\Efield_{\parallel}.
\end{equation}
% ~~~~~~~~~~~~~~~~~~~~~~~~~~~~~~~~~~~~~~~~~~~~~~~~~~~~~~~~
Це означає, що на рух, описаний у~попередньому пункті, накладається прискорений рух у~напрямку $\Efield_{\parallel} \parallel \vect{B}$. Дрейф зі швидкістю $\vect{u} = c\Efield \times \vect{B}/B^2$ зберігається, так само як і~обертання з~циклотронною частотою $\omega_c = qB/mc$.


%% --------------------------------------------------------
\section{Застосування схрещених полів}\hypertarget{crossed-fields-applications}{}
%% --------------------------------------------------------

Рух у~схрещених полях має два важливі застосування: ефект Холла в~твердому тілі та холлівські двигуни в~космічній техніці.


%% --------------------------------------------------------
\subsection*{Ефект Холла}
%% --------------------------------------------------------


Класичний приклад руху заряджених частинок у~схрещених полях реалізується у~твердому провіднику чи напівпровіднику, по якому тече електричний струм і~який знаходиться в~магнітному полі,~--- це \emphx{ефект Холла}, відкритий Е.~Холлом у~1879~р.

Розглянемо провідну пластинку завтовшки $d$ і~завширшки $w$, вздовж якої (вісь $x$) тече струм густиною $\vect{j}$, а~перпендикулярно площині пластинки (вісь $z$) прикладене магнітне поле $\vect{B}$ (рис.~\ref{tikz:hall_effect}). Носії заряду~--- концентрацією $n$ і~зарядом $q$~--- рухаються вздовж провідника з~усередненою (дрейфовою) швидкістю $\vect{v}_d$, яку підтримує поздовжнє електричне поле $\Efield_x$ (у~провіднику носії зазнають зіткнень), так що
% ~~~~~~~~~~~~~~~~~~~~~~~~~~~~~~~~~~~~~~~~~~~~~~~~~~~~~~~~
\begin{equation}\label{eq:hall_current_density}
	\vect{j} = nq\vect{v}_d.
\end{equation}
% ~~~~~~~~~~~~~~~~~~~~~~~~~~~~~~~~~~~~~~~~~~~~~~~~~~~~~~~~

\begin{figure}[h!]
	\centering
	\localinput{halleffect.tikz}
%    \includegraphics[width=0.5\linewidth]{HallEffect}
	\caption{Ефект Холла: струм $\vect{j}$ уздовж пластинки, магнітне поле $\vect{B}$ перпендикулярно їй; сила Лоренца відхиляє носії заряду до одного з~країв, у~результаті чого встановлюється поперечне холлівське поле $\Efield_H$.}
	\label{tikz:hall_effect}
\end{figure}

На кожен носій заряду, що рухається з~дрейфовою швидкістю $\vect{v}_d$ в~магнітному полі, діє сила Лоренца $\frac{q}{c}\vect{v}_d \times \vect{B}$, яка відхиляє носії до одного з~бічних країв пластинки. Внаслідок цього на протилежних гранях накопичується нескомпенсований заряд, що створює поперечне (холлівське) електричне поле $\Efield_H$. У~стаціонарному режимі результуюча поперечна сила, що діє на носії, дорівнює нулю, тобто виконується та сама умова \eqref{eq:u_condition}, яку ми використали для знаходження дрейфової швидкості:
% ~~~~~~~~~~~~~~~~~~~~~~~~~~~~~~~~~~~~~~~~~~~~~~~~~~~~~~~~
\begin{equation}\label{eq:hall_balance}
	q\Efield_H + \frac{q}{c}\vect{v}_d \times \vect{B} = 0
	\qquad \Longrightarrow \qquad
	\Efield_H = -\frac{1}{c}\vect{v}_d \times \vect{B}.
\end{equation}
% ~~~~~~~~~~~~~~~~~~~~~~~~~~~~~~~~~~~~~~~~~~~~~~~~~~~~~~~~
Отже, поперечне поле $\Efield_H$ визначається тією самою умовою нульової результуючої сили, що й~дрейф у~схрещених полях \eqref{eq:drift_velocity}. Різниця в~тому, що поздовжній рух носіїв тут підтримується поздовжнім полем $\Efield_x$ (зі зіткненнями), а~поперечне поле виникає самоузгоджено. У~сильних полях ($\omega_c\tau \gg 1$, де $\tau$~--- час між зіткненнями) аналогія з~$\vect{E}\times\vect{B}$-дрейфом буквальна; у~слабких ($\omega_c\tau \ll 1$) рівність \eqref{eq:hall_balance} виражає баланс усереднених сил.

Виразивши дрейфову швидкість через густину струму за формулою \eqref{eq:hall_current_density}, $v_d = j/(nq)$, знаходимо поперечну складову холлівського поля (тут $j = j_x$, $B = B_z$; знак $E_H$ збігається зі знаком $q$):
% ~~~~~~~~~~~~~~~~~~~~~~~~~~~~~~~~~~~~~~~~~~~~~~~~~~~~~~~~
\begin{equation}\label{eq:hall_field}
	E_H = \frac{1}{nqc}\,jB.
\end{equation}
% ~~~~~~~~~~~~~~~~~~~~~~~~~~~~~~~~~~~~~~~~~~~~~~~~~~~~~~~~
Різниця потенціалів між бічними гранями пластинки (холлівська напруга) дорівнює $U_H = E_H w$, а~оскільки повний струм $I = jwd$, отримуємо
% ~~~~~~~~~~~~~~~~~~~~~~~~~~~~~~~~~~~~~~~~~~~~~~~~~~~~~~~~
\begin{equation}\label{eq:hall_voltage}
	\tcbhighmath{
		U_H = \frac{1}{nqc}\,\frac{IB}{d}.
	}
\end{equation}
% ~~~~~~~~~~~~~~~~~~~~~~~~~~~~~~~~~~~~~~~~~~~~~~~~~~~~~~~~
Коефіцієнт
% ~~~~~~~~~~~~~~~~~~~~~~~~~~~~~~~~~~~~~~~~~~~~~~~~~~~~~~~~
\begin{equation}\label{eq:hall_coefficient}
	R_H = \frac{1}{nqc}
\end{equation}
% ~~~~~~~~~~~~~~~~~~~~~~~~~~~~~~~~~~~~~~~~~~~~~~~~~~~~~~~~
називається \emphx{сталою Холла}; вона знакова: $R_H < 0$ для електронів і~$R_H > 0$ для дірок. Ця формула справедлива для одного типу носіїв заряду з~однаковим часом релаксації; для напівпровідників з~електронами й~дірками вираз для $R_H$ складніший.

Знак холлівської напруги визначається знаком носіїв заряду: він протилежний для провідників з~електронною ($q<0$) і~дірковою ($q>0$) провідністю, що дає змогу експериментально встановити тип основних носіїв у~зразку. За відомими $I$, $B$, $d$ і~виміряним $U_H$ формула \eqref{eq:hall_voltage} дає змогу визначити концентрацію носіїв заряду $n$, а~вимірювання ефекту Холла широко використовується для характеризації напівпровідників, а~також у~датчиках магнітного поля (датчиках Холла).


%% --------------------------------------------------------
\subsection*{Дрейф у~плазмі та плазмові двигуни}
%% --------------------------------------------------------


Розглянутий електричний дрейф має фундаментальне значення для фізики плазми: оскільки швидкість дрейфу не залежить від знаку заряду й~маси частинки, замагнічені електрони та іони в~схрещених полях дрейфують \emphx{в~один і~той самий бік} з~однаковою швидкістю. Це принципово відрізняє $\vect{E}\times\vect{B}$-дрейф від руху частинки під дією самого лише електричного поля, який розділяє заряди різного знаку і~породжує струм.

Схрещені поля лежать в~основі роботи \emphx{холлівських двигунів} (\emphx{Hall-effect thrusters}), одного з~найпоширеніших типів електричних ракетних двигунів для космічних апаратів. Проте тут істотною є не однаковість дрейфу електронів та іонів, а, навпаки, різниця в~ступені їхнього замагнічення: електрони замагнічені, а~іони~--- ні.


%% --------------------------------------------------------
\paragraph{Принцип дії холлівського двигуна.}
%% --------------------------------------------------------


У~кільцевому каналі двигуна створюється радіальне магнітне поле $\vect{B}$ та осьове електричне поле $\Efield$, причому $\Efield \perp \vect{B}$. Електрони, інжектовані з~катода, потрапляючи в~область сильного радіального поля, замагнічуються (їхній ларморівський радіус $R$, див.~\eqref{eq:cyclotron_radius}, значно менший за розміри каналу) і, згідно з~формулою \eqref{eq:drift_velocity}, здійснюють азимутальний дрейф~--- це і~є \enquote{\emphx{холлівський струм}}, що дав назву двигуну.

\begin{SCfigure}[1.2][h!]
	\centering
%	\localinput{halleffectthruster.tikz}
    \includegraphics[height=8cm]{Hall_Truster}
	\caption{Схема холлівського двигуна.}
	\label{tikz:hall_thruster}
\end{SCfigure}

Іони ж, маючи значно більший ларморівський радіус (через велику масу), практично не замагнічуються в~тому самому полі й~вільно прискорюються електричним полем уздовж осі каналу, створюючи тягу. Радіальне магнітне поле, замагнічуючи електрони, різко обмежує їхню рухливість упоперек силових ліній, тобто вздовж осі до анода: замагнічений електрон обертається навколо силової лінії і~зміщується впоперек неї лише внаслідок зіткнень. Тому основне падіння потенціалу зосереджується в~зоні прискорення поблизу виходу каналу. Електрони, що дрейфують, одночасно іонізують ударами робочу речовину (зазвичай ксенон), а~плазма в~каналі в~цілому залишається квазінейтральною.

%\paragraph{Магнітне утримання та дрейфи в термоядерній плазмі.}
%Аналогічний $\vect{E}\times\vect{B}$-дрейф відіграє ключову роль в установках магнітного утримання плазми (токамаки, стеларатори). Тут, окрім електричного дрейфу, суттєвими є й інші дрейфи того самого типу --- \emphx{градієнтний дрейф} (при неоднорідності $\vect{B}$) і \emphx{дрейф кривини} магнітних силових ліній, які отримуються аналогічним чином, якщо в рівнянні \eqref{eq:u_condition} замінити електричну силу $q\Efield$ на відповідну ефективну силу:
%\begin{equation}\label{eq:general_drift}
%	\vect{u}_{\vect{F}} = \frac{c}{q}\,\frac{\vect{F} \times \vect{B}}{B^2}.
%\end{equation}
%Зокрема, для градієнтного дрейфу $\vect{F} = -\mu\nabla B$ (де $\mu = mv_\perp^2/2B$ --- магнітний момент частинки), а для дрейфу кривини $\vect{F} = mv_\parallel^2 \vect{R}_c/R_c^2$, де $\vect{R}_c$ --- радіус-вектор кривини силової лінії, напрямлений від центра кривини до точки, де розглядається частинка. Оскільки ці дрейфи, на відміну від електричного, \emphx{залежать від знаку заряду}, вони спричиняють поляризацію плазми та є однією з причин турбулентного переносу, що обмежує якість утримання в термоядерних реакторах.
%
%Таким чином, простий приклад руху зарядженої частинки в схрещених полях виявляється не лише навчальною задачею класичної електродинаміки, а й основою для розуміння як інженерних застосувань (електричні ракетні двигуни), так і фундаментальних проблем керованого термоядерного синтезу.


%% --------------------------------------------------------
\section{Методи вимірювання питомого заряду електрона}\hypertarget{charge-to-mass-ratio}{}
%% --------------------------------------------------------

\emphx{Питомим зарядом} називається відношення $q/m$, де $m$~--- маса частинки.

Рух частинки за наявності тільки електричних і~магнітних полів описується рівнянням
% ~~~~~~~~~~~~~~~~~~~~~~~~~~~~~~~~~~~~~~~~~~~~~~~~~~~~~~~~
\begin{equation}\label{eq:q_m_motion}
	m\frac{d\vect{v}}{dt} = q\left(\Efield + \frac{1}{c}\vect{v} \times \vect{B}\right).
\end{equation}
% ~~~~~~~~~~~~~~~~~~~~~~~~~~~~~~~~~~~~~~~~~~~~~~~~~~~~~~~~
Звідси видно, що воно повністю визначається питомим зарядом $q/m$, а~не зарядом і~масою окремо.


%% --------------------------------------------------------
\subsection*{Вимірювання $q/m$ за відхиленням траєкторії в~електричному полі}
%% --------------------------------------------------------


У~цьому методі пучок частинок пропускається між обкладками зарядженого конденсатора, і~за виміряним відхиленням траєкторії визначають відношення $q/m$. Нехай $l$~--- довжина конденсатора вздовж початкової швидкості частинок $\vect{v}_0$, а~$L$~--- відстань від краю конденсатора до екрана; швидкість $v_0$ вважаємо відомою. Схема досліду показана на рис.~\ref{tikz:condenser_deflection}.

\begin{figure}[h!]
	\centering
	\localinput{In_condensator.tikz}
	\caption{Відхилення зарядженої частинки електричним полем конденсатора.}
	\label{tikz:condenser_deflection}
\end{figure}

У~конденсаторі є однорідне електричне поле з~напруженістю $\Efield$, напрямлене перпендикулярно початковій швидкості частинки $\vect{v}_0$. Це поле викликає рух з~прискоренням
% ~~~~~~~~~~~~~~~~~~~~~~~~~~~~~~~~~~~~~~~~~~~~~~~~~~~~~~~~
\begin{equation}\label{eq:a_E}
	\vect{a} = \frac{q}{m}\Efield, \qquad \vect{a} \perp \vect{v}_0.
\end{equation}
% ~~~~~~~~~~~~~~~~~~~~~~~~~~~~~~~~~~~~~~~~~~~~~~~~~~~~~~~~
Якби поле було відсутнє, то частинки потрапляли б у~точку \enquote{0} екрана Е (рис.~\ref{tikz:condenser_deflection}). Під дією поля відхилення частинки в~конденсаторі у~вертикальному напрямку становить
% ~~~~~~~~~~~~~~~~~~~~~~~~~~~~~~~~~~~~~~~~~~~~~~~~~~~~~~~~
\begin{equation}\label{eq:y1_E}
	y_1 = \frac{at^2}{2} = \frac{qE}{2m}t^2 = \frac{q}{2m}E\frac{l^2}{v_0^2},
\end{equation}
% ~~~~~~~~~~~~~~~~~~~~~~~~~~~~~~~~~~~~~~~~~~~~~~~~~~~~~~~~
де $t = l/v_0$~--- час прольоту конденсатора. Оскільки поле $\Efield$ напрямлене вздовж осі $y$, компонента швидкості $v_x = \const = v_0$.

На виході з~конденсатора швидкість $v_y$ становитиме
% ~~~~~~~~~~~~~~~~~~~~~~~~~~~~~~~~~~~~~~~~~~~~~~~~~~~~~~~~
\begin{equation}\label{eq:vy_E}
	v_y = at = \frac{qE}{m}\frac{l}{v_0}.
\end{equation}
% ~~~~~~~~~~~~~~~~~~~~~~~~~~~~~~~~~~~~~~~~~~~~~~~~~~~~~~~~
Далі частинка летить прямолінійно, причому траєкторія утворює кут $\alpha$ з~віссю $x$ такий, що
% ~~~~~~~~~~~~~~~~~~~~~~~~~~~~~~~~~~~~~~~~~~~~~~~~~~~~~~~~
\begin{equation}\label{eq:tg_alpha}
	\tg\alpha = \frac{v_y}{v_x} = \frac{qE}{m}\frac{l}{v_0^2}; \qquad
	y_1 = \frac{l}{2}\tg\alpha.
\end{equation}
% ~~~~~~~~~~~~~~~~~~~~~~~~~~~~~~~~~~~~~~~~~~~~~~~~~~~~~~~~

Долітаючи до екрана, вона набуває додаткового відхилення
% ~~~~~~~~~~~~~~~~~~~~~~~~~~~~~~~~~~~~~~~~~~~~~~~~~~~~~~~~
\begin{equation}\label{eq:y2_E}
	y_2 = L \tg\alpha.
\end{equation}
% ~~~~~~~~~~~~~~~~~~~~~~~~~~~~~~~~~~~~~~~~~~~~~~~~~~~~~~~~
Повне зміщення частинки в~результаті становитиме
% ~~~~~~~~~~~~~~~~~~~~~~~~~~~~~~~~~~~~~~~~~~~~~~~~~~~~~~~~
\begin{equation}\label{eq:y_total_E}
	\tcbhighmath{
		y = y_1 + y_2 = \left(\frac{1}{2}l + L\right)\tg\alpha = \frac{qE}{mv_0^{2}}l\left(\frac{1}{2}l + L\right).
	}
\end{equation}
% ~~~~~~~~~~~~~~~~~~~~~~~~~~~~~~~~~~~~~~~~~~~~~~~~~~~~~~~~
Таким чином, за виміряним відхиленням частинки $y$ за допомогою формули \eqref{eq:y_total_E} визначається питомий заряд частинки $q/m$.

Швидкість $v_0$ тут має бути відома окремо. Якщо, наприклад, частинки прискорено напругою $U$ ($mv_0^2/2 = qU$), то $q/m$ у~формулі \eqref{eq:y_total_E} скорочується: $y = \dfrac{El(l/2 + L)}{2U}$, тобто відхилення в~електричному полі від питомого заряду не залежить.


%% --------------------------------------------------------
\subsection*{Вимірювання $q/m$ за відхиленням траєкторії в~магнітному полі}
%% --------------------------------------------------------


Для вимірювання питомого заряду можна замість електричного поля використати магнітне, також напрямлене перпендикулярно траєкторії частинки (рис.~\ref{tikz:magnetic_deflection}). На частинку діє сила Лоренца $\vect{F} = \frac{q}{c}\vect{v} \times \vect{B}$, створюючи прискорення
% ~~~~~~~~~~~~~~~~~~~~~~~~~~~~~~~~~~~~~~~~~~~~~~~~~~~~~~~~
\begin{equation}\label{eq:a_B}
	\vect{a} = \frac{\vect{F}}{m} = \frac{q}{mc}\vect{v} \times \vect{B}.
\end{equation}
% ~~~~~~~~~~~~~~~~~~~~~~~~~~~~~~~~~~~~~~~~~~~~~~~~~~~~~~~~

\begin{figure}[h!]
	\centering
	\localinput{In_magfield.tikz}
	\caption{Відхилення зарядженої частинки магнітним полем.}
	\label{tikz:magnetic_deflection}
\end{figure}

Нехай швидкість частинки достатньо велика, так що радіус кривини траєкторії $R = mcv_0/(|q|B)$ набагато більший за довжину $l$ області поля, а~відхилення за час прольоту мале. Тоді можна прийняти
% ~~~~~~~~~~~~~~~~~~~~~~~~~~~~~~~~~~~~~~~~~~~~~~~~~~~~~~~~
\begin{equation}\label{eq:a_B_approx}
	v_x \approx v_0, \qquad
	\vect{a} \approx \frac{q}{mc}\vect{v}_0 \times \vect{B}, \qquad
	a = |\vect{a}| \approx \frac{|q|}{mc}v_0 B.
\end{equation}
% ~~~~~~~~~~~~~~~~~~~~~~~~~~~~~~~~~~~~~~~~~~~~~~~~~~~~~~~~
Відхилення частинки за час прольоту поля $t = l/v_0$ становитиме
% ~~~~~~~~~~~~~~~~~~~~~~~~~~~~~~~~~~~~~~~~~~~~~~~~~~~~~~~~
\begin{equation}\label{eq:y1_B}
	y_1 = \frac{1}{2}at^2 = \frac{q}{2m}B\frac{l^2}{v_0 c}.
\end{equation}
% ~~~~~~~~~~~~~~~~~~~~~~~~~~~~~~~~~~~~~~~~~~~~~~~~~~~~~~~~
Швидкість $v_y$ частинки на виході з~області дії поля дорівнює
% ~~~~~~~~~~~~~~~~~~~~~~~~~~~~~~~~~~~~~~~~~~~~~~~~~~~~~~~~
\begin{equation}\label{eq:vy_B}
	v_y = at = \frac{ql}{mc}B,
\end{equation}
% ~~~~~~~~~~~~~~~~~~~~~~~~~~~~~~~~~~~~~~~~~~~~~~~~~~~~~~~~
а~кут, який утворює вектор швидкості з~віссю $x$, визначається рівністю
% ~~~~~~~~~~~~~~~~~~~~~~~~~~~~~~~~~~~~~~~~~~~~~~~~~~~~~~~~
\begin{equation}\label{eq:tg_beta}
	\tg\beta = \frac{v_y}{v_x} \approx \frac{v_y}{v_0} = \frac{qlB}{mv_0 c}, \qquad
	y_1 = \frac{l}{2}\tg\beta.
\end{equation}
% ~~~~~~~~~~~~~~~~~~~~~~~~~~~~~~~~~~~~~~~~~~~~~~~~~~~~~~~~
Долітаючи до екрана, частинка отримує додаткове зміщення
% ~~~~~~~~~~~~~~~~~~~~~~~~~~~~~~~~~~~~~~~~~~~~~~~~~~~~~~~~
\begin{equation}\label{eq:y2_B}
	y_2 = L \tg\beta.
\end{equation}
% ~~~~~~~~~~~~~~~~~~~~~~~~~~~~~~~~~~~~~~~~~~~~~~~~~~~~~~~~
Повне зміщення частинки становитиме
% ~~~~~~~~~~~~~~~~~~~~~~~~~~~~~~~~~~~~~~~~~~~~~~~~~~~~~~~~
\begin{equation}\label{eq:y_total_B}
	\tcbhighmath{
		y = y_1 + y_2 = \left(\frac{1}{2}l + L\right)\tg\beta = \frac{qB}{mv_0 c}l\left(\frac{1}{2}l + L\right).
	}
\end{equation}
% ~~~~~~~~~~~~~~~~~~~~~~~~~~~~~~~~~~~~~~~~~~~~~~~~~~~~~~~~
Вимірюючи це зміщення, можна знайти питомий заряд $q/m$ частинки.

Тут теж потрібна відома швидкість $v_0$. Якщо частинки прискорено напругою $U$, то $v_0 = \sqrt{2qU/m}$ і~$y = \sqrt{q/m}\,\dfrac{Bl}{c\sqrt{2U}}\left(\dfrac{l}{2} + L\right)$, тобто магнітне відхилення, на відміну від електричного, залежить від $q/m$.


%% --------------------------------------------------------
\subsection*{Метод Томсона}
%% --------------------------------------------------------


Дж.~Дж.~Томсон використав компенсаційну схему, в~якій рух частинки визначався як електричним, так і~магнітним полем. Ця схема давала змогу не вимірювати попередньо швидкість частинок $v_0$. Суть методу полягає в~такому: спочатку вмикають тільки магнітне поле (як описано вище) і~вимірюють відхилення пучка заряджених частинок:
% ~~~~~~~~~~~~~~~~~~~~~~~~~~~~~~~~~~~~~~~~~~~~~~~~~~~~~~~~
\begin{equation}\label{eq:Thomson_B}
	y = \frac{qB}{mv_0 c}l\left(\frac{1}{2}l + L\right).
\end{equation}
% ~~~~~~~~~~~~~~~~~~~~~~~~~~~~~~~~~~~~~~~~~~~~~~~~~~~~~~~~
Потім додатково вмикають електричне поле, перпендикулярне як до магнітного поля ($\Efield \perp \vect{B}$), так і~до швидкості пучка, і~підбирають його величину таким чином, щоб повністю компенсувати відхилення пучка і~повернути його в~точку $y = 0$. Умовою компенсації є рівність
% ~~~~~~~~~~~~~~~~~~~~~~~~~~~~~~~~~~~~~~~~~~~~~~~~~~~~~~~~
\begin{equation}\label{eq:Thomson_compensation}
	qE = \frac{q}{c}v_0 B, \qquad \text{або} \qquad v_0 = c\frac{E}{B}.
\end{equation}
% ~~~~~~~~~~~~~~~~~~~~~~~~~~~~~~~~~~~~~~~~~~~~~~~~~~~~~~~~
Тому наведений вище вираз для зміщення пучка в~магнітному полі переписується у~вигляді
% ~~~~~~~~~~~~~~~~~~~~~~~~~~~~~~~~~~~~~~~~~~~~~~~~~~~~~~~~
\begin{equation}\label{eq:Thomson_y}
	y = \frac{q}{m}\frac{B^2}{c^2 E}l\left(\frac{1}{2}l + L\right).
\end{equation}
% ~~~~~~~~~~~~~~~~~~~~~~~~~~~~~~~~~~~~~~~~~~~~~~~~~~~~~~~~
Звідси за відомими значеннями зміщення $y$ і~полів $B$ та $E$ визначається питомий заряд $q/m$.

Умова \eqref{eq:Thomson_compensation} збігається з~умовою \eqref{eq:u_condition}: частинки, що не відхиляються, рухаються з~дрейфовою швидкістю \eqref{eq:drift_velocity}, тобто $v_0 = u$. Схрещені поля відбирають частинки за швидкістю, тому таку конфігурацію називають \emphx{фільтром швидкостей}.

Описаним методом Дж.~Дж.~Томсон у~1897~р. при вивченні катодних променів вперше виміряв питомий заряд електрона. Сучасне значення питомого заряду електрона (Томсон отримав величину того ж порядку; тут $e < 0$~--- заряд електрона) становить
% ~~~~~~~~~~~~~~~~~~~~~~~~~~~~~~~~~~~~~~~~~~~~~~~~~~~~~~~~
\begin{equation}\label{eq:e_m_value}
	e/m = -5{,}273 \cdot 10^{17} \ \text{од.~СГСЕ} \cdot \text{г}^{-1}
	= -1{,}759 \cdot 10^{11} \ \text{Кл} \cdot \text{кг}^{-1}.
\end{equation}
% ~~~~~~~~~~~~~~~~~~~~~~~~~~~~~~~~~~~~~~~~~~~~~~~~~~~~~~~~


%% --------------------------------------------------------
\section{Досліди Міллікена}\hypertarget{millikan-experiment}{}
%% --------------------------------------------------------

У~попередніх дослідах вимірювався питомий заряд частинок. Р.~Міллікен у~1909~р. запропонував спосіб безпосереднього вимірювання заряду. Схема його установки показана на рис.~\ref{tikz:Millikan}.

У~простір між пластинами конденсатора Міллікен вводив за допомогою розпилювача найдрібніші крапельки олії. Краплі електризувалися при розпиленні, а~їхній заряд можна було змінювати, опромінюючи повітря в~камері рентгенівськими променями: утворені іони захоплювалися краплями. Рухом крапель керували, змінюючи величину і~знак напруги на пластинах конденсатора.

\begin{figure}[h!]
	\centering
	\includegraphics[width=0.75\linewidth]{Milliken}
	\caption{Схема установки Міллікена для вимірювання заряду електрона (сила $q\Efield$ вказана для випадку $q < 0$).}
	\label{tikz:Millikan}
\end{figure}

Розглянемо метод докладніше. Будемо для визначеності вважати, що заряд краплі від'ємний, $q < 0$. У~формулах нижче, де йдеться про величини сил і~швидкостей, під $q$ розуміємо $|q|$.

На краплю діють такі сили:
\begin{enumerate}
	\item сила з~боку електричного поля $q\Efield$;
	\item рівнодійна сили тяжіння і~архімедової сили
	% ~~~~~~~~~~~~~~~~~~~~~~~~~~~~~~~~~~~~~~~~~~~~~~~~~~~~~~~~
	\begin{equation}\label{eq:gravity_force}
		\vect{F}_{g} = \frac{4\pi}{3}r^3(\rho - \rho_0)\vect{g},
	\end{equation}
	% ~~~~~~~~~~~~~~~~~~~~~~~~~~~~~~~~~~~~~~~~~~~~~~~~~~~~~~~~
	де $r$~--- радіус краплі, $\rho$ і~$\rho_0$~--- густини краплі і~повітря відповідно;
	\item сила в'язкого тертя
	% ~~~~~~~~~~~~~~~~~~~~~~~~~~~~~~~~~~~~~~~~~~~~~~~~~~~~~~~~
	\begin{equation}\label{eq:viscous_force}
		\vect{F}_{\eta} = -6\pi\eta r\vect{v}.
	\end{equation}
	% ~~~~~~~~~~~~~~~~~~~~~~~~~~~~~~~~~~~~~~~~~~~~~~~~~~~~~~~~
\end{enumerate}

За відсутності електричного поля крапля падає з~постійною швидкістю $v_0$, яка визначається з~умови $F_g = F_\eta$, або
% ~~~~~~~~~~~~~~~~~~~~~~~~~~~~~~~~~~~~~~~~~~~~~~~~~~~~~~~~
\begin{equation}\label{eq:falling_equilibrium}
	\frac{4\pi}{3}r^3(\rho - \rho_0)g = 6\pi\eta r v_0.
\end{equation}
% ~~~~~~~~~~~~~~~~~~~~~~~~~~~~~~~~~~~~~~~~~~~~~~~~~~~~~~~~

Увімкнувши електричне поле, змушують краплю підніматися з~невеликою швидкістю $v_E$, яку можна знайти з~рівняння
% ~~~~~~~~~~~~~~~~~~~~~~~~~~~~~~~~~~~~~~~~~~~~~~~~~~~~~~~~
\begin{equation}\label{eq:rising_equilibrium}
	\frac{4\pi}{3}r^3(\rho - \rho_0)g + 6\pi\eta r v_E = qE.
\end{equation}
% ~~~~~~~~~~~~~~~~~~~~~~~~~~~~~~~~~~~~~~~~~~~~~~~~~~~~~~~~
З~урахуванням рівності \eqref{eq:falling_equilibrium} перепишемо це співвідношення у~вигляді
% ~~~~~~~~~~~~~~~~~~~~~~~~~~~~~~~~~~~~~~~~~~~~~~~~~~~~~~~~
\begin{equation}\label{eq:charge_balance}
	6\pi\eta r(v_0 + v_E) = qE,
\end{equation}
% ~~~~~~~~~~~~~~~~~~~~~~~~~~~~~~~~~~~~~~~~~~~~~~~~~~~~~~~~
звідки можна знайти заряд краплі:
% ~~~~~~~~~~~~~~~~~~~~~~~~~~~~~~~~~~~~~~~~~~~~~~~~~~~~~~~~
\begin{equation}\label{eq:charge_drop}
	q = 6\pi\eta r\frac{v_0 + v_E}{E}.
\end{equation}
% ~~~~~~~~~~~~~~~~~~~~~~~~~~~~~~~~~~~~~~~~~~~~~~~~~~~~~~~~

Практично вимірювання здійснювалося таким способом. За відсутності електричного поля вимірювалася швидкість вільного падіння краплі $v_0$. Потім із співвідношення \eqref{eq:falling_equilibrium} визначався радіус краплі:
% ~~~~~~~~~~~~~~~~~~~~~~~~~~~~~~~~~~~~~~~~~~~~~~~~~~~~~~~~
\begin{equation}\label{eq:drop_radius}
	r = \sqrt{\frac{9\eta v_0}{2g(\rho - \rho_0)}}.
\end{equation}
% ~~~~~~~~~~~~~~~~~~~~~~~~~~~~~~~~~~~~~~~~~~~~~~~~~~~~~~~~
Увімкнувши електричне поле і~вимірявши швидкість підйому краплі $v_E$, за формулами \eqref{eq:charge_drop}--\eqref{eq:drop_radius} або
% ~~~~~~~~~~~~~~~~~~~~~~~~~~~~~~~~~~~~~~~~~~~~~~~~~~~~~~~~
\begin{equation}\label{eq:charge_full}
	q = 6\pi\eta\sqrt{\frac{9\eta v_0}{2g(\rho - \rho_0)}}\frac{v_0 + v_E}{E}
	= 9\pi\sqrt{\frac{2\eta^3 v_0}{g(\rho - \rho_0)}}\frac{v_0 + v_E}{E}
\end{equation}
% ~~~~~~~~~~~~~~~~~~~~~~~~~~~~~~~~~~~~~~~~~~~~~~~~~~~~~~~~
можна знайти заряд $q$. Слід мати на увазі, що для крапель мікронних розмірів закон Стокса потребує поправки (Каннінгема) на проковзування газу, оскільки радіус краплі порівнянний з~довжиною вільного пробігу молекул повітря; Міллікен враховував таку поправку в~обробленні результатів.

Заряд краплі змінювався при захопленні іона, і~при цьому стрибком змінювалася швидкість підйому $v_E$. Згідно з~\eqref{eq:charge_drop} стрибку $\Delta v_E$ відповідає зміна заряду $\Delta q = 6\pi\eta r\,\Delta v_E/E$. Аналізуючи результати численних вимірювань, Міллікен встановив, що зміни заряду крапель $\Delta q$ і~сам заряд $q$ були завжди цілими кратними одній і~тій же мінімальній величині:
% ~~~~~~~~~~~~~~~~~~~~~~~~~~~~~~~~~~~~~~~~~~~~~~~~~~~~~~~~
\begin{equation}\label{eq:elementary_charge}
	|e| \approx 4{,}80 \cdot 10^{-10} \ \text{Фр} = 1{,}602 \cdot 10^{-19} \ \text{Кл}.
\end{equation}
% ~~~~~~~~~~~~~~~~~~~~~~~~~~~~~~~~~~~~~~~~~~~~~~~~~~~~~~~~
За відомими значеннями питомого заряду $e/m$ і~заряду $e$ електрона можна знайти масу:
% ~~~~~~~~~~~~~~~~~~~~~~~~~~~~~~~~~~~~~~~~~~~~~~~~~~~~~~~~
\begin{equation}\label{eq:electron_mass}
	\frac{|e|}{|e/m|} = m \approx 9{,}11 \cdot 10^{-28} \ \text{г}.
\end{equation}
% ~~~~~~~~~~~~~~~~~~~~~~~~~~~~~~~~~~~~~~~~~~~~~~~~~~~~~~~~
Значення досліду Міллікена полягає в~тому, що була експериментально доведена дискретність електричного заряду.

\begin{Historical}
Дослід виконано Р.~Міллікеном разом із Г.~Флетчером у~1909--1913~рр. Спершу використовували краплі води, але вони швидко випаровувалися, тому перейшли до олії, краплі якої тривалий час зберігають розмір. Значення заряду, отримане Міллікеном у~1913~р. ($4{,}774\cdot10^{-10}$~од.~СГСЕ), було занижене приблизно на $0{,}6\,\%$ через неточне значення в'язкості повітря. Результати Міллікена оскаржував Ф.~Еренгафт, який стверджував, що спостерігав \enquote{субелектрони} з~меншим зарядом; подальші вимірювання цього не підтвердили. У~1923~р. Міллікен отримав Нобелівську премію з~фізики.
\end{Historical}


%% --------------------------------------------------------
\section*{Підсумки}


\phantomsection
\addcontentsline{toc}{section}{Підсумки}
%% --------------------------------------------------------

Рух зарядженої частинки в~електричному й~магнітному полях визначається її питомим зарядом $q/m$ та початковими умовами: саме ця комбінація, а~не заряд і~маса окремо, входить у~рівняння руху. Основні результати:
\begin{itemize}
	\item однорідне $\Efield$~--- рівноприскорений рух \eqref{eq:motion_E_solution};
	\item однорідне $\vect{B}$~--- рівномірний рух уздовж поля й~обертання з~циклотронною частотою $|\omega_c| = |q|B/mc$ по колу ларморівського радіуса $R = mcv_{\perp}/(|q|B)$, тобто гвинтова лінія; сила Лоренца не виконує роботи;
	\item $\Efield \parallel \vect{B}$~--- спіраль зі зростаючим кроком;
	\item $\Efield \perp \vect{B}$, $E < B$~--- до обертання додається дрейф $\vect{u} = c\,\Efield \times \vect{B}/B^2$, що не залежить ні від знака заряду, ні від маси; траєкторія~--- трохоїда, а~за $|v_0| = |u|$ (зокрема з~нульовою початковою швидкістю)~--- циклоїда;
	\item ефект Холла: $U_H = \dfrac{1}{nqc}\dfrac{IB}{d}$, $R_H = \dfrac{1}{nqc}$; знак $U_H$ визначає тип носіїв заряду;
	\item метод Томсона: компенсація відхилення дає $v_0 = cE/B$ (фільтр швидкостей) і~питомий заряд $q/m$; дослід Міллікена: $q = 6\pi\eta r\,(v_0 + v_E)/E$, заряд краплі кратний $e$.
\end{itemize}
Ці два вимірювання дають змогу знайти й~саму масу електрона: $m = |e|/|e/m| \approx 9{,}11\cdot10^{-28}$~г.


%% --------------------------------------------------------
\section*{Задачі}


\phantomsection
\addcontentsline{toc}{section}{Задачі}
%% --------------------------------------------------------

\begin{problem}\label{prob:ElectronInField}
	\emphx{Електрон в~однорідному полі.} Електрон влітає зі швидкістю $v=10^{9}$~см/с в~однорідне магнітне поле $B=100$~Гс перпендикулярно до ліній індукції. (а)~Покажіть, що траєкторія~--- коло, і~знайдіть його радіус $R$ та циклотронну частоту $\omega_c$; чи залежить $\omega_c$ від $v$? (б)~Обчисліть $R$ і~$\omega_c$ ($e=4{,}8\cdot10^{-10}$~СГСЕ, $m=9{,}1\cdot10^{-28}$~г, $c=3\cdot10^{10}$~см/с). (в)~Яким буде рух, якщо $\vect{v}$ утворює з~$\vect{B}$ кут $\alpha$?
	% \emphx{Відповідь:} (а)~$R=\dfrac{mvc}{eB}$, $\omega_c=\dfrac{v}{R}=\dfrac{eB}{mc}$ --- від $v$ не залежить. (б)~$R\approx0{,}57$~см, $\omega_c\approx1{,}8\cdot10^{9}$~с$^{-1}$ ($T\approx3{,}6$~нс). (в)~Гвинтова лінія з радіусом $R\sin\alpha$ і кроком $h=2\pi R\cos\alpha$.
\end{problem}

\begin{problem}\label{prob:HelixEB}
	\emphx{Спіраль у~паралельних полях.} Частинка із зарядом $q$ і~масою $m$ починає рух зі швидкістю $v_0$, перпендикулярною до однорідних паралельних полів $\Efield \parallel \vect{B}$. Знайдіть її траєкторію та зміщення $\Delta z_n$ вздовж поля за $n$-й оберт.
	% \emphx{Відповідь:} коло радіуса $R = \dfrac{mcv_0}{|q|B}$ в площині, перпендикулярній до $\vect{B}$, і рівноприскорений рух вздовж поля: $z = \dfrac{qE}{2m}t^2$. За період $T = \dfrac{2\pi mc}{|q|B}$ зміщення $\Delta z_n = (2n-1)\,\dfrac{2\pi^2 mc^2E}{|q|B^2}$ (у напрямку сили $q\Efield$), тобто зміщення за послідовні оберти відносяться як $1:3:5:\dots$
\end{problem}

\begin{problem}\label{prob:CycloidElectron}
	\emphx{Електрон у~схрещених полях.} Електрон стартує зі стану спокою в~схрещених полях $E = 0{,}1$~СГСЕ ($30$~В/см) і~$B = 100$~Гс. Знайдіть (а)~швидкість і~напрямок дрейфу; (б)~максимальну швидкість електрона; (в)~максимальне відхилення від осі дрейфу; (г)~зсув вздовж осі дрейфу за один період обертання. В~який бік спершу відхиляється електрон?
	% \emphx{Відповідь:} (а)~$u = cE/B \approx 3\cdot10^{7}$~см/с у напрямку $\Efield \times \vect{B}$. (б)~$v_{\max} = 2u \approx 6\cdot10^{7}$~см/с. (в)~$2u/|\omega_c| \approx 0{,}034$~см (радіус \enquote{колеса} $u/|\omega_c| \approx 0{,}017$~см). (г)~$2\pi u/|\omega_c| = uT \approx 0{,}11$~см при $T \approx 3{,}6$~нс. Електрон спершу рухається в бік $-\Efield$ (сила $q\Efield$ для $q<0$), а дрейфує в напрямку $\Efield \times \vect{B}$ так само, як і позитивний заряд.
\end{problem}

\begin{problem}\label{prob:HallSensor}
	\emphx{Датчик Холла.} Напівпровідникова пластинка з~електронною провідністю має концентрацію носіїв $n = 10^{16}$~см$^{-3}$ і~товщину $d = 0{,}5$~мм. Крізь неї тече струм $I = 1$~мА, а~перпендикулярно до пластинки прикладено поле $B = 1000$~Гс. (а)~Обчисліть холлівську напругу. (б)~Як змінить знак $U_H$, якщо замінити пластинку діркового типу з~тією самою концентрацією? (в)~Як за виміряною $U_H$ знайти $n$? (Використайте $1$~мА $= 3\cdot10^{6}$~СГСЕ, $1$~СГСЕ потенціалу $= 300$~В.)
	\emphx{Відповідь:} (а)~$U_H = \dfrac{IB}{n|q|cd} \approx 4{,}2\cdot10^{-7}$~СГСЕ $\approx 0{,}13$~мВ. (б)~Знак змінюється на протилежний ($R_H>0$ для дірок). (в)~$n = \dfrac{IB}{|q|cd\,|U_H|}$.
\end{problem}

\begin{problem}\label{prob:HallThrusterMagnetization}
	\emphx{Замагніченість у~холлівському двигуні.} У~каналі двигуна радіальне поле $B = 200$~Гс, осьове електричне поле $E = 200$~В/см. Оцініть (а)~ларморівський радіус електрона з~енергією $10$~еВ і~іона ксенону ($M \approx 131$~а.о.м.) з~енергією $300$~еВ; (б)~швидкість азимутального дрейфу електронів. Чи можна вважати іони замагніченими, якщо розмір каналу становить кілька сантиметрів?
	\emphx{Відповідь:} (а)~$R_e \approx 0{,}5$~мм, $R_i \approx 1{,}4$~м. (б)~$u = cE/B \approx 10^{8}$~см/с ($E = 0{,}667$~СГСЕ). Оскільки $R_i \gg$ розмірів каналу, іони практично не замагнічені, а~електрони ($R_e \ll$ розмірів каналу) замагнічені.
\end{problem}

\begin{problem}\label{prob:DeflectionAccelerated}
	\emphx{Відхилення після прискорення.} Електрони і~протони прискорено однаковою напругою $U$ і~пропущено крізь відхиляючий конденсатор, а~в~другому досліді~--- крізь область магнітного поля (довжини $l$, відстань до екрана $L$, відхилення мале). Покажіть, що (а)~в електричному полі відхилення $y$ не залежить від $q/m$, і~знайдіть $y$; (б)~у магнітному полі $|y_e|/|y_p| = \sqrt{m_p/m_e}$.
	% \emphx{Відповідь:} (а)~З $mv_0^2/2 = |q|U$ маємо $y = \dfrac{El(l/2+L)}{2U}$ (частинки різних знаків відхиляються в протилежні боки). (б)~$|y| = \sqrt{|q|/m}\,\dfrac{Bl}{c\sqrt{2U}}\left(\dfrac{l}{2}+L\right)$, тому $|y_e|/|y_p| = \sqrt{m_p/m_e} \approx 43$.
\end{problem}

\begin{problem}\label{prob:ThomsonNumbers}
	\emphx{Метод Томсона.} У~схемі Томсона довжина області полів $l = 2$~см, відстань до екрана $L = 20$~см. При $B = 3$~Гс пляма на екрані зміщується на $y = 2{,}2$~см, а~електричне поле $E = 0{,}1$~СГСЕ повертає її в~початкове положення. Знайдіть швидкість електронів і~питомий заряд.
	\emphx{Відповідь:} $v_0 = cE/B \approx 10^{9}$~см/с; $\left|\dfrac{q}{m}\right| = \dfrac{y\,v_0\,c}{B\,l\,(l/2+L)} \approx 5{,}2\cdot10^{17}$~од.~СГСЕ$\cdot$г$^{-1}$ $\approx 1{,}7\cdot10^{11}$~Кл$\cdot$кг$^{-1}$.
\end{problem}

\begin{problem}\label{prob:OilDrop}
	\emphx{Краплина Міллікена.} Крапля олії ($\rho = 0{,}9$~г/см$^3$) у~повітрі ($\eta = 1{,}8\cdot10^{-4}$~г$\cdot$см$^{-1}\cdot$с$^{-1}$, густиною повітря знехтуйте) падає зі швидкістю $v_0 = 0{,}02$~см/с. У~полі $E = 10$~СГСЕ вона піднімається зі швидкістю $v_E = 0{,}011$~см/с, а~після захоплення іона~--- зі швидкістю $v_E' = 0{,}0214$~см/с. Знайдіть радіус краплі, її заряд у~першому випадку та зміну заряду. Поправкою Каннінгема знехтуйте.
	% \emphx{Відповідь:} $r = \sqrt{\dfrac{9\eta v_0}{2g\rho}} \approx 1{,}4\cdot10^{-4}$~см; $q = 6\pi\eta r\,\dfrac{v_0+v_E}{E} \approx 1{,}4\cdot10^{-9}$~СГСЕ $\approx 3e$; $\Delta q = 6\pi\eta r\,\dfrac{v_E'-v_E}{E} \approx 4{,}8\cdot10^{-10}$~СГСЕ $\approx e$.
\end{problem}

\begin{Questions}
\begin{quizlist}
	\item Запишіть рівняння руху зарядженої частинки в~однорідному електричному полі та його розв'язок. Яким буде рух і~куди напрямлене прискорення?
	\item Чому сила Лоренца не виконує роботи і~не змінює модуль швидкості частинки? Як розкладають швидкість на складові $\vect{v}_{\parallel}$ і~$\vect{v}_{\perp}$ в~магнітному полі та який рух відповідає кожній із них?
	\item Що таке циклотронна частота? Від яких величин вона залежить, а~від яких (наприклад, від швидкості частинки)~--- ні? Запишіть формулу для радіуса кола $R$.
	\item Якою є траєкторія частинки в~однорідних паралельних полях $\Efield \parallel \vect{B}$ (рис.~\ref{tikz:helix})? Чому крок спіралі зростає?
	\item У~схрещених полях $\Efield \perp \vect{B}$ введено швидкість $\vect{u}$ (рис.~\ref{tikz:drift}). З~якої умови її знаходять, і~чому їй відповідає формула $\vect{u} = c\,\Efield\times\vect{B}/B^2$? Куди напрямлений дрейф відносно $\Efield$ і~$\vect{B}$?
	\item Чому швидкість електричного дрейфу не залежить від знака заряду й~маси частинки? Чому при цьому вона все ж не поділяє заряди різного знака, на відміну від руху під дією самого лише електричного поля?
	\item З~яких трьох рухів складається рух у~схрещених полях? За якої умови існує дрейфова швидкість, що відбувається при $E > B$ і~коли потрібні релятивістські рівняння?
	\item Що таке трохоїда та циклоїда? За якої початкової умови траєкторія в~схрещених полях є циклоїдою (рис.~\ref{tikz:drift_cycloide}) і~з~чим її можна порівняти механічно?
	\item У~чому полягає ефект Холла (рис.~\ref{tikz:hall_effect})? Чому в~стаціонарному режимі результуюча поперечна сила на носії дорівнює нулю?
	\item Виведіть формулу для холлівської напруги $U_H$ і~сталої Холла $R_H$. Які величини можна визначити за виміряною $U_H$ і~як за знаком холлівської напруги визначити тип носіїв?
	\item Поясніть принцип дії холлівського двигуна. Чому електрони в~каналі замагнічені, а~іони~--- ні, і~як це забезпечує тягу?
	\item Що таке питомий заряд? Чому рух частинки в~полях визначається саме відношенням $q/m$, а~не зарядом і~масою окремо?
	\item Опишіть метод вимірювання $q/m$ за відхиленням в~електричному полі конденсатора (рис.~\ref{tikz:condenser_deflection}). З~яких двох частин складається повне зміщення $y = y_1 + y_2$?
	\item У~чому перевага методу Томсона порівняно з~вимірюванням окремо в~електричному та магнітному полі? Яка умова компенсації відхилення, що з~неї випливає для швидкості $v_0$ і~чому таку конфігурацію полів називають фільтром швидкостей?
	\item Опишіть установку Міллікена (рис.~\ref{tikz:Millikan}). Які сили діють на краплю і~як за швидкостями $v_0$ та $v_E$ визначають її радіус і~заряд?
	\item Що довели досліди Міллікена і~чому саме зміни заряду крапель були кратні одній і~тій же величині? За відомими $e$ та $e/m$ як знайти масу електрона?
\end{quizlist}
\end{Questions}